\section{Introduction}

Language-model agents are now asked to judge whether a research idea is new: AI-scientist systems check the novelty of the ideas they generate before executing them \citep{lu2024the,yamada2025the}, and a growing set of tools assesses ideas or submissions against the literature \citep{shahid2025literaturegrounded,moussa2025scholareval,zhang2026opennovelty,hou2026memonoveltyagent,qiao2026innoeval}. The failure these tools must avoid is concrete and has been observed: an independent evaluation of the AI Scientist found that ``several generated research ideas were incorrectly classified as novel, including well-established concepts such as micro-batching for stochastic gradient descent'' \citep{beel2025evaluating}, and a survey of 24 runnable AI-scientist systems found that only 38\% report any novelty-verification method \citep{ding2026autonomous}.

Two observations motivate this paper. First, evaluations of novelty judgement rest on human labels, which are hard to obtain and contested: LLM judges and domain experts reach opposite conclusions about the novelty of generated research questions \citep{sinhahajari2026on}, LLM judges are swayed by style \citep{ji2026style}, and even experts find novelty hard to judge \citep{si2024can}. Second, a large diagnostic study of research agents attributes their failures to a missing \emph{metacognitive loop}---``the ability to check what they produced against what they found, revise when it does not hold up''---places the deficit at the model level, and explicitly leaves open ``whether orchestration-level interventions can close it'' \citep{fei2026how}. Novelty assessors in the literature are, by their own descriptions, fixed retrieve-compare-report pipelines; we found none that searches for evidence against its \emph{own} verdict, or that measures what doing so changes (\S\ref{sec:related}).

We address both with a checkable task and a controlled comparison.

\paragraph{A temporal prior-art task.} Whether an idea has \emph{already been published} has an objective answer when the realising paper and its date are known. We take recent papers, rewrite each paper's core idea as a proposal that names nothing identifying, and run an agent on the proposal twice: once with the literature visible up to today (\textsc{scooped}: the right outcome is to find the realising paper and conclude the idea is explored) and once with every work published on or after the realising paper's date hidden at the literature layer (\textsc{prepub}: the paper cannot be found, so an ``already explored'' verdict must rest on earlier work). We call this Temporal Prior-Art Detection (TPAD).

\paragraph{A controlled ablation of self-challenge.} We instrument an open research-investigation harness, ResearchForge, whose controller can (i) raise its own questions that challenge the critique, each gap and each new direction with a contradiction-seeking literature search, and (ii) revise its critique when such a search weakens a conclusion. Switches turn these off, or replace the adaptive controller with a fixed pipeline, while holding the model, tools, evidence rules and budgets fixed.

\paragraph{Contributions.}
\begin{enumerate}
  \item The TPAD protocol, with anti-leakage controls (a publication-date cutoff enforced at every literature API and on every stored record, first-version full text, disabled code search, a closed-book probe of the agent's model) and a scorer that reads only the agent's structured records (\S\ref{sec:benchmark}).
  \item An instrumented harness in which self-challenge, critique revision, adaptive control and source-integrity checks are independent switches (\S\ref{sec:system}).
  \item A pilot study with NVIDIA Nemotron~3 Super on \PilotItems{} items from papers published in September 2026 (\S\ref{sec:results}). We report it as a pilot: effect estimates with uncertainty, failure analysis, and what a full study needs.
\end{enumerate}
We release the harness, items, run traces and scoring code.
